\documentclass[11pt]{article}
\usepackage{amsmath,amssymb,amsthm,hyperref}
\usepackage[margin=1in]{geometry}
\newtheorem{theorem}{Theorem}
\DeclareMathOperator{\tr}{tr}
\title{Scalar counterexamples to the signed polynomial trace inequality}
\author{Principia Math In Progress}
\date{29 September 2026}
\begin{document}\maketitle
\begin{abstract}
The polynomial clause of the signed mean value trace conjecture of
Paulin, Mackey, and Tropp fails for every integer degree $q\ge2$, even
for real scalar matrices. At degree three it would assert $7\le3$.
This refutes the joint conjecture recorded as MathDB 325507.
The separate exponential clause is not resolved here.
\end{abstract}

\section*{Task 1: counterexample to the joint conjecture}
For a Hermitian matrix $X$, let $X_+$ and $X_-$ be its positive and
negative parts, so $X=X_+-X_-$ and $|X|=X_++X_-$. Put $D=A-B$.
Conjecture 7.3 of \cite{PMT13}, page 17, and Conjecture 13.4 of
\cite{PMT16}, pages 3465--3466, assert two inequalities for arbitrary
Hermitian $A,B,C$, every integer $q\ge1$, and every $s>0$.
The polynomial clause is
\begin{equation}\label{P}
\tr[C(A^q-B^q)]\le\frac q2\tr\!\left[
 (sD_+^2+s^{-1}C_+^2)|A|^{q-1}
 +(sD_-^2+s^{-1}C_-^2)|B|^{q-1}\right].
\end{equation}
The other clause replaces $A^q,B^q$ by $e^A,e^B$, replaces the two
absolute-power weights by $e^A,e^B$, and replaces $q/2$ by $1/2$.
Neither source restricts $A$ or $B$ to be positive semidefinite.
The final article separately asserts validity for scalars; the examples
below also contradict that assertion for the polynomial clause.

\begin{theorem}
Inequality \eqref{P} holds universally when $q=1$, but fails for each
integer $q\ge2$. In particular, $A=-2$, $B=-1$, $C=-1$, $q=3$, $s=1$
is a scalar counterexample.
\end{theorem}
\begin{proof}
For this explicit example, $D=-1$, $D_+=C_+=0$, and $D_-=C_-=1$.
The left side of \eqref{P} is
\[
 (-1)((-2)^3-(-1)^3)=7,
\]
whereas the right side is $(3/2)(1+1)\cdot1=3$.
All five inputs satisfy the hypotheses, and $7>3$.

The same choice $A=-2$, $B=-1$, $C=-1$, $s=1$ gives left side
$2^q-1$ and right side $q$ for every odd $q\ge3$.
The binomial theorem gives
$2^q=1+q+\sum_{j=2}^q\binom qj>1+q$, so each such degree fails.

For even $q\ge2$, take instead
\[
 A=-4,\qquad B=-1,\qquad C=1,\qquad s=\frac{2^{q-1}}3.
\]
Now $D_+=C_-=0$, $D_-=3$, and $C_+=1$. The two sides are
\[
 L=4^q-1,\qquad
 R=\frac q2\left(s^{-1}4^{q-1}+9s\right)=3q2^{q-1}.
\]
The $q$ positive numbers $1,4,\ldots,4^{q-1}$ are not all equal.
The strict arithmetic--geometric mean inequality therefore yields
\[
 4^q-1=3\sum_{j=0}^{q-1}4^j
 >3q\left(\prod_{j=0}^{q-1}4^j\right)^{1/q}
 =3q2^{q-1}=R.
\]
This proves failure for every even degree.

Finally, for $q=1$ the weights are identity matrices and
$D_+^2+D_-^2=D^2$, $C_+^2+C_-^2=C^2$. Since the trace of the square
of a Hermitian matrix is nonnegative,
\[
 0\le\tr\!\left[(\sqrt{s}D-s^{-1/2}C)^2\right]
 =s\tr D^2+s^{-1}\tr C^2-2\tr(CD),
\]
which is precisely \eqref{P}.
\end{proof}

These scalar examples can also be multiplied by the identity matrix
of any fixed positive dimension; both sides are then multiplied by
that dimension. Thus the obstruction does not depend on permitting
dimension one. Failure of the polynomial clause is enough to refute
the conjunction of the two displayed universal inequalities. It does
not refute the exponential inequality by itself, and no conclusion
about its validity is claimed here. The unsigned inequalities proved
elsewhere in the source are different statements.

\section*{Four proposed follow-up tasks}
\begin{enumerate}
\item[Task 2.] Prove or refute the exponential clause for arbitrary
Hermitian $A,B,C$ and all $s>0$.
\item[Task 3.] Determine whether the polynomial clause holds when
$A$ and $B$ are positive semidefinite, with arbitrary Hermitian $C$.
\item[Task 4.] Determine which differentiable real functions $f$ obey
the corresponding scalar signed inequality with weights $f'(a)$ and
$f'(b)$ for all real $a,b,c$ and $s>0$.
\item[Task 5.] Find a valid signed replacement for the polynomial
clause that yields one-sided matrix moment bounds in the intended
exchangeable-pair application.
\end{enumerate}

\begin{thebibliography}{9}
\bibitem{PMT13} D.~Paulin, L.~Mackey, and J.~A.~Tropp,
\emph{Deriving Matrix Concentration Inequalities from Kernel Couplings},
\href{https://arxiv.org/abs/1305.0612v1}{arXiv:1305.0612v1} (2013),
Conjecture 7.3, p.~17.
\bibitem{PMT16} D.~Paulin, L.~Mackey, and J.~A.~Tropp,
\emph{Efron--Stein inequalities for random matrices},
Ann. Probab. \textbf{44} (2016), 3431--3473,
\href{https://doi.org/10.1214/15-AOP1054}{doi:10.1214/15-AOP1054},
Conjecture 13.4, pp.~3465--3466.
\end{thebibliography}
\end{document}
